\section*{第 \ref{ch:b2:diffcalc} 章——微分学}

\begin{solution}{exo:b2:diffcalc:1}
$J_f = \begin{pmatrix} 2x & -2y\\ 2y & 2x\end{pmatrix}$, $\det J_f
= 4(x^2 + y^2)$：除原点外可逆。（此处 $f$ 为 $z
\mapsto z^2$ 的复数形式。）

$J_\Phi = \begin{pmatrix} \cos\theta & -r\sin\theta & 0\\
\sin\theta & r\cos\theta & 0\\ 0 & 0 & 1\end{pmatrix}$, $\det = r$:
当 $r \neq 0$ 时可逆（柱坐标）。
\end{solution}

\begin{solution}{exo:b2:diffcalc:2}
对 $v = (a, b) \neq 0$：$\frac{f(tv) - 0}{t} = \frac{t^3a^3}{t\cdot
t^2(a^2+b^2)} = \frac{a^3}{a^2 + b^2}$：各方向导数均存在，
且其值为 $D_v = \frac{a^3}{a^2+b^2}$。但 $v
\mapsto D_v$ 不是线性的（$D_{(1,0)} = 1$，$D_{(0,1)} = 0$，
$D_{(1,1)} = \frac12 \neq 1$）：任何线性映射都无法给出这些
值，故 $f$ 在 $0$ 处不\omterm{def:b2:diffcalc:differential}{可微}（\omterm{def:b2:diffcalc:differential}{微分}
必须是 $v \mapsto D_v$）。
\end{solution}

\begin{solution}{exo:b2:diffcalc:3}
$f = x^3 + y^3 - 3xy$：临界点为 $(0,0)$ 与 $(1,1)$（本科一年级
的计算）。\omterm{thm:b2:diffcalc:taylor}{黑塞矩阵}为 $H = \begin{pmatrix} 6x & -3\\ -3 &
6y\end{pmatrix}$。在 $(0,0)$ 处\omterm{def:b2:reduction:eigen}{特征值}异号（$\det = -9 < 0$）：
鞍点。在 $(1,1)$ 处，$\det = 27 > 0$，迹 $> 0$：正
定，严格局部极小值——现由
\cref{thm:b2:diffcalc:extrema}证明，而非直接断定。

$g = x^4 + y^4 - 2(x-y)^2$: $\nabla g = (4x^3 - 4(x - y),\; 4y^3 +
4(x-y))$；临界点为 $(0,0)$、$(\sqrt2, -\sqrt2)$、
$(-\sqrt2, \sqrt2)$（本科一年级的计算）。在 $(\pm\sqrt2, \mp\sqrt2)$ 处，$H =
\begin{pmatrix} 12\cdot2 - 4 & 4\\ 4 & 20\end{pmatrix} =
\begin{pmatrix} 20 & 4\\ 4 & 20\end{pmatrix}$：正定
（对角占优；\omterm{def:b2:reduction:eigen}{特征值}为 $24, 16$）：严格局部极小值。
在 $(0,0)$ 处，$H = \begin{pmatrix} -4 & 4\\ 4 & -4\end{pmatrix}$，
奇异、半负定：判别法失效；方向
研究（$g(x,x) = 2x^4 > 0$，$g(x,-x) = 2x^4 - 8x^2 < 0$，取值充分小时）
表明其为类鞍点——不是极值点。
\end{solution}

\begin{solution}{exo:b2:diffcalc:4}
对 $t \mapsto f(tx)$ 在 $t = 1$ 处求导：由链式法则，
$\langle \nabla f(x), x\rangle$；由齐次性，同一函数
即 $t^pf(x)$，其导数为 $pf(x)$（取 $t = 1$）：欧拉恒等式。

反之：固定 $x \neq 0$，令 $\varphi(t) = f(tx) - t^p f(x)$
在 $t > 0$ 上，则 $\varphi'(t) = \langle\nabla f(tx), x\rangle -
pt^{p-1}f(x) = \frac1t\bigl(\langle \nabla f(tx), tx\rangle -
p\,t^pf(x)\bigr)$。由假设，即在点
$tx$ 处的欧拉恒等式，将括号化为 $p\,f(tx) - p\,t^pf(x) =
p\,\varphi(t)$。于是 $\varphi' = \frac{p}{t}\varphi$，且 $\varphi(1)
= 0$：线性常\omterm{def:b2:diffcalc:differential}{微分}方程的唯一解为 $\varphi \equiv 0$（本科一年级
的唯一性），即 $f(tx) = t^pf(x)$。
\end{solution}

\begin{solution}{exo:b2:diffcalc:5}
展开 $f(x + h) - f(x) = \langle Ax - b, h\rangle +
\frac12\langle Ah, h\rangle$（由 $A$ 的对称性）：$\nabla f(x) = Ax -
b$，且处处有 $H_f = A$。$f$ 为凸函数当且仅当 $A$ 是半
正定（\omterm{thm:b2:diffcalc:taylor}{黑塞矩阵}判别法，因 $H$ 为常值，故此处全局成立：
二阶泰勒公式精确成立）。若 $A$ 为
正定，则唯一临界点为 $x^* = A^{-1}b$，且 $f(x^*
+ h) - f(x^*) = \frac12\langle Ah, h\rangle \geq
\frac{\lambda_{\min}}{2}\norm h^2 > 0$（$h \neq 0$）：严格
全局极小值。
\end{solution}

\begin{solution}{exo:b2:diffcalc:6}
由于 $\nabla g(a) \neq 0$，不妨设其中一个偏导数 $\frac{\partial
g}{\partial y}(a) \neq 0$。由隐函数定理
（与 \cref{thm:b2:diffcalc:inverse} 配套）可得 $\{g =
0\}$ 在 $a = (a_1, a_2)$ 附近可写为 $\gamma(t) = (t, y(t))$，其中 $y$
为 $C^1$ 类，$y'(t) = -\frac{\partial_x g}{\partial_y g}(\gamma(t))$
（对 $g(t, y(t)) = 0$ 求导）。一元函数 $t
\mapsto f(\gamma(t))$ 在 $t = a_1$ 处取局部极值：
\[
0 = \frac{\dd}{\dd t}f(\gamma(t))\Big|_{a_1}
= \partial_x f(a) + \partial_y f(a)\,y'(a_1)
= \partial_x f(a) - \partial_yf(a)\frac{\partial_x
g(a)}{\partial_y g(a)} :
\]
由此，向量 $\nabla f(a)$ 与 $\nabla g(a)$ 的
分量成比例：$\nabla f(a) = \lambda \nabla g(a)$，其中 $\lambda =
\frac{\partial_y f(a)}{\partial_y g(a)}$。

应用：在圆周上，$\nabla(xy) = (y, x)$ 与
$(2x, 2y)$ 平行，推出 $y^2 = x^2$；结合约束条件，候选点为
$\pm\bigl(\tfrac{1}{\sqrt2}, \tfrac{1}{\sqrt2}\bigr)$（函数值
$\frac12$）及 $\pm\bigl(\tfrac{1}{\sqrt2},
-\tfrac{1}{\sqrt2}\bigr)$（函数值 $-\frac12$）：最大值 $\frac12$，最小
值 $-\frac12$（可取到，因圆周为紧集）。
\end{solution}

\begin{solution}{exo:b2:diffcalc:7}
$\nabla f = (2x - y + 1,\; 2y - x - 1) = 0$：解得 $x =
-\frac13$，$y = \frac13$。\omterm{thm:b2:diffcalc:taylor}{黑塞矩阵}为 $\begin{pmatrix} 2 & -1\\ -1 &
2\end{pmatrix}$，正定：二次函数
$f$ 的全局极小值，值为 $f\bigl(-\frac13, \frac13\bigr) = -\frac13$。

在三角形 $T$ 上：内部临界点 $(-\frac13,
\frac13) \notin T$（$x$ 为负）。各边：在 $y = 0$ 上，$x \in
\intcc{0}{1}$：$f = x^2 + x$，单调递增，端点值为 $0$ 与 $2$。在
$x = 0$ 上：$f = y^2 - y$，最小值 $-\frac14$ 在 $y = \frac12$ 处取得，
两端点处的值为 $0$ 与 $0$。在 $x + y = 1$ 上：代入 $y = 1 - x$，
$f = x^2 + (1-x)^2 - x(1-x) + x - (1-x) = 3x^2 - x$；在
$\intcc{0}{1}$ 上：最小值 $-\frac{1}{12}$ 在 $x = \frac16$ 处取得，
端点处的值分别为 $0$（$x=0$ 时）与 $2$（$x=1$ 时）。顶点：$f(0,0) = 0$，$f(1,0)
= 2$，$f(0,1) = 0$。在 $T$ 上：全局最小值 $-\frac14$ 在 $(0,
\frac12)$ 处取得，最大值 $2$ 在 $(1, 0)$ 处。
\end{solution}

\begin{solution}{exo:b2:diffcalc:8}
$J_f = \begin{pmatrix} 1 & 2y\\ 2x & 1\end{pmatrix}$, $\det J_f =
1 - 4xy$。在 $0$ 处，$\det = 1 \neq 0$，为局部\omterm{def:b2:diffcalc:differential}{微分}同胚
（由\cref{thm:b2:diffcalc:inverse}），且
\[
\dd(f^{-1})_{(0,0)} = (J_f(0))^{-1} = I_2 .
\]
球上的可逆性：需 $4\abs{xy} < 1$ 处处成立；在
$\norm{(x,y)}_2 < r$ 上，$\abs{xy} \leq \frac{x^2 + y^2}{2} <
\frac{r^2}{2}$，故 $r = \frac{1}{\sqrt2}$ 可取；且它是最大的：
在 $(x, y) = \bigl(\tfrac12, \tfrac12\bigr)$ 处，\omterm{def:b2:nvs:norm}{范数}为
$\frac{1}{\sqrt2}$，$\det J_f = 0$。
\end{solution}

\begin{solution}{exo:b2:diffcalc:9}
$f(t) = (\cos t, \sin t)$：$f(0) = f(2\pi) = (1, 0)$，然而 $f'(t) =
(-\sin t, \cos t)$ 的\omterm{def:b2:nvs:norm}{范数}为 $1$，从不为零：不存在使
导数为零的点——罗尔定理没有向量类比。均值
不等式显然成立：$\norm{f(2\pi) - f(0)} = 0 \leq
2\pi \cdot \sup\norm{f'} = 2\pi$。
\end{solution}

\begin{solution}{exo:b2:diffcalc:10}
$f(x+h) - f(x) = 2\langle x, h\rangle + \norm h^2$: $\dd f_x =
2\langle x, \cdot\rangle$，$\nabla f(x) = 2x$，\omterm{thm:b2:diffcalc:taylor}{黑塞矩阵} $2I$
（常值）。对 $g$：
\[
g(x + h) - g(x) = \langle Ax, h\rangle + \langle Ah, x\rangle
+ \langle Ah, h\rangle
= \bigl\langle (A + A^{\mathsf T})x,\ h\bigr\rangle +
O(\norm h^2),
\]
故 $\nabla g(x) = (A + A^{\mathsf T})x$，且 $H_g = A +
A^{\mathsf T}$，为常值。由\cref{exo:b2:diffcalc:11}，$g$ 为
凸函数当且仅当 $A + A^{\mathsf T}$ 半正定——
只有 $A$ 的对称部分起作用，正如 $g(x) =
\langle \frac{A + A^{\mathsf T}}2 x, x\rangle$。
\end{solution}

\begin{solution}{exo:b2:diffcalc:11}
$f$ 为凸函数当且仅当其在每条线段上的限制为凸函数，
即每个 $\varphi(t) = f(a + tv)$ 为凸函数。由链式
法则得 $\varphi''(t) = \langle H_{a+tv}\,v,\ v\rangle$。

若所有\omterm{thm:b2:diffcalc:taylor}{黑塞矩阵}都半正定：$\varphi'' \geq 0$，
故每个 $\varphi$ 为凸函数（本科一年级分册），从而 $f$ 为凸函数。
反之，若 $f$ 为凸函数，则每个 $\varphi$ 为凸，故
$\varphi''(0) \geq 0$：$\langle H_a v, v\rangle \geq 0$ 对
每个 $a$ 与每个方向 $v$，所有\omterm{thm:b2:diffcalc:taylor}{黑塞矩阵}
都半正定。
\end{solution}

\begin{solution}{exo:b2:diffcalc:12}
\begin{enumerate}
  \item 由均值不等式
        （\cref{thm:b2:diffcalc:mvi}）应用于 $g$：
        $\norm{g(x) - g(y)} \leq k\norm{x-y}$，故
        \[
        \norm{f(x) - f(y)} \geq \norm{x - y} - \norm{g(x) -
        g(y)} \geq (1 - k)\norm{x - y} :
        \]
        $f$ 为单射，且 $f^{-1}$（定义于像集上）
        为 $\frac{1}{1-k}$-利普希茨映射。
  \item 固定 $y$，$T(x) = y - g(x)$ 满足 $\norm{T(x) -
        T(x')} = \norm{g(x') - g(x)} \leq k\norm{x - x'}$：一个
        $\R^n$ 这一\omterm{def:b2:metric:complete}{完备}空间上的压缩映射。巴拿赫
        （\cref{thm:b2:metric:banach}）给出一个不动点 $x^*
        = y - g(x^*)$，即 $f(x^*) = y$：$f$ 为满射。
  \item 于是 $f$ 为 $\R^n$ 上的双射，其
        逆为 $\frac{1}{1-k}$-利普希茨映射：一个整体反函数
        定理，其中 $\dd g$ 处处足够小，
        取代了
        \cref{thm:b2:diffcalc:inverse} 中的局部可逆性假设。
\end{enumerate}
\end{solution}

\begin{solution}{pb:b2:diffcalc:1}
\textbf{1.} $\norm{ABx} \leq \vertiii A\norm{Bx} \leq
\vertiii A\vertiii B\norm x$：对 $\norm x = 1$ 取上确界。
矩阵乘积、$\det$ 与 $\operatorname{tr}$ 都是
元素的多项式函数，从而\omterm{def:b2:metric:continuity}{连续}（$\mathcal M_n(\R)
\simeq \R^{n^2}$，所有\omterm{thm:b2:nvs:finitedim}{范数等价}
（\cref{thm:b2:nvs:finitedim}）。

\textbf{2.} $\sum\vertiii{X^k} \leq \sum\vertiii X^k <
\infty$：级数\omterm{def:b2:series:def}{绝对收敛}，从而收敛
（\cref{thm:b2:nvs:absoluteconvergence}）。由 $(I -
X)\sum_{k\leq N}X^k = I - X^{N+1} \to I$：和为 $(I -
X)^{-1}$。\omterm{def:b2:nvs:norm}{范数}：$\leq \sum\vertiii X^k = \frac{1}{1 -
\vertiii X}$；且
\[
(I - X)^{-1} - I - X = \sum_{k\geq2}X^k = X^2(I - X)^{-1},
\qquad
\vertiii{X^2(I-X)^{-1}} \leq
\frac{\vertiii X^2}{1 - \vertiii X} = O(\vertiii X^2).
\]

\textbf{3.} $A + H = A(I + A^{-1}H)$，其中 $\vertiii{A^{-1}H}
\leq \vertiii{A^{-1}}\vertiii H < 1$：由第 2 问知其可逆：
半径为 $\vertiii{A^{-1}}^{-1}$、以 $A$ 为心的开球位于
$GL_n(\R)$ 中。稠密性：$\det(A + \varepsilon I)$ 是 $n$ 次
多项式（关于 $\varepsilon$），首项系数为 $1$：它
只有有限多个根，故存在 $\varepsilon_k \to 0$
使 $A + \varepsilon_kI$ 可逆，且收敛于 $A$。

\textbf{4.} 对 $\vertiii H < \frac{1}{2\vertiii{A^{-1}}}$：
\[
(A + H)^{-1} - A^{-1}
= \bigl[(I + A^{-1}H)^{-1} - I\bigr]A^{-1},
\]
其\omterm{def:b2:nvs:norm}{范数}至多为 $\frac{\vertiii{A^{-1}H}}{1 -
\vertiii{A^{-1}H}}\,\vertiii{A^{-1}} \leq
2\vertiii{A^{-1}}^2\vertiii H \to 0$：$\Phi$ 在
每个 $A \in GL_n(\R)$ 处\omterm{def:b2:metric:continuity}{连续}。

\textbf{5.} $(A+H)^2 = A^2 + AH + HA + H^2$：映射 $H
\mapsto AH + HA$ 是线性的，误差 $H^2$ 为
$O(\vertiii H^2)$。展开 $(A + H)^k$，按 $H$ 因子的个数归类，线性部分为
$\sum_{i=0}^{k-1}A^iHA^{k-1-i}$，含 $\geq 2$ 个 $H$ 因子的项
以 $\binom k2$ 个
\omterm{def:b2:nvs:norm}{范数} $\leq \vertiii A^{k-2}\vertiii H^2$ 量级的乘积为界，
即 $O(\vertiii H^2)$。
不能把和式合并为 $kA^{k-1}H$，因为 $H$ 与 $A$
未必可交换——这个和式才是正确的非交换
导数。

\textbf{6.} 对充分小的 $H$：
\[
(A+H)^{-1} = (I + A^{-1}H)^{-1}A^{-1}
= \bigl(I - A^{-1}H + O(\vertiii H^2)\bigr)A^{-1}
= A^{-1} - A^{-1}HA^{-1} + O(\vertiii H^2) :
\]
$\dd\Phi_A(H) = -A^{-1}HA^{-1}$，关于 $H$ 线性。当 $n = 1$ 时：
$\dd(1/a)(h) = -h/a^2$，即熟悉的导数。

\textbf{7.} 沿曲线的链式法则：
$\bigl(A(t)^{-1}\bigr)' = \dd\Phi_{A(t)}(A'(t)) =
-A(t)^{-1}A'(t)A(t)^{-1}$。当 $A(t) = I + tB$，$t = 0$ 时：$(I +
tB)^{-1} = I - tB + O(t^2)$。

\textbf{8.} 由第 5 题及迹的循环不变性：
\[
\dd f_A(H) = \operatorname{tr}\Bigl(\sum_{i=0}^{k-1}
A^iHA^{k-1-i}\Bigr) = k\operatorname{tr}\bigl(A^{k-1}H\bigr) .
\]
对照弗罗贝尼乌斯内积，$\dd f_A(H) =
\operatorname{tr}\bigl((\nabla f)^{\mathsf T}H\bigr)$ 要求
$(\nabla f)^{\mathsf T} = kA^{k-1}$，于是 $\nabla f(A) =
k\,(A^{k-1})^{\mathsf T}$。

\textbf{9.} $f(A + H) - f(A) = 2\operatorname{tr}
(A^{\mathsf T}H) + \operatorname{tr}(H^{\mathsf T}H)$：其
\omterm{def:b2:diffcalc:differential}{微分}为 $H \mapsto 2\operatorname{tr}(A^{\mathsf T}H) =
2\langle A, H\rangle$，故 $\nabla f(A) = 2A$；二阶项
恰为 $\norm H_F^2$：\omterm{thm:b2:diffcalc:taylor}{黑塞矩阵}
为恒同\omterm{def:b2:quadratic:def}{二次型}的两倍，正定且为常值，
因此 $\norm\cdot_F^2$ 严格凸（
泰勒公式在此恰好精确）。

\textbf{10.} 由关于列的多重线性，$\det(I + H) =
\sum_{S\subseteq\{1,\dots,n\}}\det(M_S)$，其中 $M_S$ 的列
为 $h_j$（$j \in S$），其余为 $e_j$。$S = \varnothing$
时为 $1$；$S = \{j\}$ 时，为将 $I$ 的第 $j$ 列
换成 $h_j$ 所得的\omterm{def:b2:linalg:det}{行列式}，即其第 $j$ 个元素 $h_{jj}$，
求和得 $\operatorname{tr}H$；每个满足 $\abs S \geq 2$ 的项
是至少两列大小为
$O(\vertiii H)$ 的\omterm{def:b2:linalg:det}{行列式}，因而为 $O(\vertiii H^2)$（有限维空间上的多重线性映射有界），
且这样的项有限多个。
因此 $\det(I + H) = 1 + \operatorname{tr}H +
O(\vertiii H^2)$：$\dd(\det)_I = \operatorname{tr}$。

\textbf{11.} $\det(A + H) = \det A\,\det(I + A^{-1}H) =
\det A\,\bigl(1 + \operatorname{tr}(A^{-1}H) +
O(\vertiii H^2)\bigr)$：\omterm{def:b2:diffcalc:differential}{微分}为 $H \mapsto
\det(A)\operatorname{tr}(A^{-1}H)$。

\textbf{12.} 沿第 $i$ 行的拉普拉斯展开：$\det A =
\sum_j a_{ij}C_{ij}$，且代数余子式 $C_{ij}$ 不涉及
第 $i$ 行：$\frac{\partial\det}{\partial a_{ij}} = C_{ij}$。由此
\[
\dd(\det)_A(H) = \sum_{i,j}C_{ij}h_{ij}
= \operatorname{tr}\bigl(\operatorname{com}(A)^{\mathsf T}
H\bigr)
= \operatorname{tr}\bigl(\operatorname{adj}(A)H\bigr),
\qquad
\nabla(\det)(A) = \operatorname{com}(A) .
\]
对可逆 $A$，$\operatorname{adj}A = \det(A)A^{-1}$
与第 11 题一致。沿 $C^1$ 曲线，链式法则为
$(\det A(t))' = \operatorname{tr}(\operatorname{adj}
(A(t))\,A'(t))$：\omterm{pb:b2:diffcalc:1}{雅可比公式}。

\textbf{13.} 设 $A' = MA$，且
$\operatorname{adj}(A)A = \det(A)I$：
\[
(\det A)' = \operatorname{tr}\bigl(\operatorname{adj}(A)MA
\bigr)
= \operatorname{tr}\bigl(A\operatorname{adj}(A)M\bigr)
= \det A\;\operatorname{tr}M
\]
（由迹的循环性；$A\operatorname{adj}A = \det(A) I$ 同样成立）。
标量线性常\omterm{def:b2:diffcalc:differential}{微分}方程 $y' = \operatorname{tr}(M(t))\,y$ 的
唯一解为 $y(t) = y(0)\exp\bigl(\int_0^t
\operatorname{tr}M\bigr)$（本科一年级）：刘维尔公式。

\textbf{14.} 在 $A \in SL_n(\R)$ 处取 $H = \frac1nA$：
$\dd(\det)_A\bigl(\tfrac1nA\bigr) = \frac1n\det(A)
\operatorname{tr}(A^{-1}A) = \frac1n\cdot1\cdot n = 1 \neq 0$：
该\omterm{def:b2:diffcalc:differential}{微分}是非零线性型，因而在水平集的每一点处都满射到
$\R$ 上：$SL_n(\R)$ 是一个光滑的水平集
（维数为 $n^2 - 1$）。

\textbf{15.} $\sum_k\vertiii{A^k/k!} \leq
\sum\vertiii A^k/k! = \eu^{\vertiii A}$：\omterm{def:b2:series:def}{绝对收敛}
（第 2 题的\omterm{def:b2:metric:complete}{完备}性论证），
在每个球 $\vertiii A \leq R$ 上\omterm{def:b2:funcseq:series}{正规收敛}（界 $R^k/k!$ 与
$A$ 无关）。每个部分和\omterm{def:b2:metric:continuity}{连续}（为多项式）；
\omterm{def:b2:funcseq:def}{一致收敛}的极限在各球上\omterm{def:b2:metric:continuity}{连续}：$A \mapsto \eu^A$ 是
\omterm{def:b2:metric:continuity}{连续}的，且 $\vertiii{\eu^A} \leq \eu^{\vertiii A}$。

\textbf{16.} 两级数都\omterm{def:b2:series:def}{绝对收敛}，故
可以作\omterm{thm:b2:series:fubini}{柯西乘积}（见\cref{thm:b2:series:fubini}）：
\[
\eu^A\eu^B = \sum_{n\geq0}\frac{1}{n!}\sum_{k=0}^n\binom
nkA^kB^{n-k}
= \sum_{n\geq0}\frac{(A+B)^n}{n!} = \eu^{A+B},
\]
二项式恒等式要求 $AB = BA$。取 $B = -A$：
$\eu^A\eu^{-A} = \eu^0 = I$：每个 $\eu^A \in GL_n(\R)$。

\textbf{17.} 级数 $\sum t^kA^k/k!$ 及其导出级数
$\sum t^{k-1}A^k/(k-1)! = A\sum t^{k-1}A^{k-1}/(k-1)!$ 在
每个区段 $\abs t \leq T$ 上\omterm{def:b2:funcseq:series}{正规收敛}（界为
$T^k\vertiii A^k/k!$）：级数的求导定理
（\cref{thm:b2:funcseq:seriestransfer}，逐元素应用）
给出 $\frac{\dd}{\dd t}\eu^{tA} = A\eu^{tA}$；把 $A$
提到右边则得 $\eu^{tA}A$。逐次求导得 $C^\infty$。

\textbf{18.} $A(t) = \eu^{tA}$ 满足 $A'(t) = A\,A(t)$：
由刘维尔公式（第 13 题），取常数 $M = A$，得
$\det\eu^{tA} = \eu^{t\operatorname{tr}A}$（值为 $1$，位于 $t =
0$ 处）；取 $t = 1$，$\det\eu^A = \eu^{\operatorname{tr}A}$。
检验：$n = 1$ 时即指数函数本身；幂零矩阵 $A$ 满足
$\operatorname{tr}A = 0$，且 $\eu^A$ 是\omterm{def:b2:linalg:det}{行列式}为
$1$；$\operatorname{tr}A = 0$ 给出 $\det\eu^A = 1$：
迹为零的矩阵被映到 $SL_n(\R)$ 中。

\textbf{19.} $\eu^H - I - H = \sum_{k\geq2}H^k/k!$，其\omterm{def:b2:nvs:norm}{范数}
$\leq \vertiii H^2\eu^{\vertiii H} = O(\vertiii H^2)$，
故 $\dd(\exp)_0 = \mathrm{id}$ 可逆。又 $\exp$ 是
$C^1$：由第 5 题，\omterm{def:b2:diffcalc:differential}{微分}的候选 $H \mapsto
\sum_k\frac1{k!}\sum_iA^iHA^{k-1-i}$ 是\omterm{def:b2:funcseq:series}{正规收敛}的
线性映射的级数，它关于 $A$ \omterm{def:b2:metric:continuity}{连续}依赖（在球上界为
$\vertiii A^{k-1}/(k-1)!$），因此偏导数存在且
都\omterm{def:b2:metric:continuity}{连续}（\cref{thm:b2:diffcalc:c1} 及级数
传递定理）。由反函数定理
（\cref{thm:b2:diffcalc:inverse}）在 $0$ 处适用：$\exp$ 是
一个 $C^1$ \omterm{def:b2:diffcalc:differential}{微分}同胚，将 $0$ 的某邻域
映为 $I$ 的某邻域——$I$ 附近的矩阵都有对数。

\textbf{20.} 对对称矩阵 $S = PDP^{\mathsf T}$（谱
定理）：$\eu^S = P\eu^DP^{\mathsf T}$ 是对称的，\omterm{def:b2:reduction:eigen}{特征值}为
$\eu^{\lambda_i} > 0$，故正定。
满射性：正定矩阵 $Q =
P\operatorname{diag}(\mu_i)P^{\mathsf T}$（$\mu_i > 0$）是
$\eu^S$ 的值，其中 $S = P\operatorname{diag}(\ln\mu_i)P^{\mathsf T}$。
单射性：$\eu^S$ 决定其\omterm{def:b2:reduction:eigen}{特征子空间}，
它们恰为 $S$ 的\omterm{def:b2:reduction:eigen}{特征子空间}（在 $S$ 对应 $\lambda$ 的每个\omterm{def:b2:reduction:eigen}{特征子空间}上，
$\eu^S$ 作用为 $\eu^\lambda$；不同的 $\lambda$ 给出
不同的 $\eu^\lambda$），对\omterm{def:b2:reduction:eigen}{特征值}取 $\ln$
即可还原 $S$。因此 $\exp$ 是从对称矩阵
到正定矩阵上的双射。

\textbf{21.} $F(A + H) = A^{\mathsf T}A + A^{\mathsf T}H +
H^{\mathsf T}A + H^{\mathsf T}H$: $\dd F_A(H) =
A^{\mathsf T}H + H^{\mathsf T}A$（取值于 $S_n$；误差
$O(\vertiii H^2)$）。在 $A \in O_n$ 处，对 $S \in S_n$，
取 $H = \frac12AS$，则
\[
A^{\mathsf T}\cdot\tfrac12AS + \tfrac12(AS)^{\mathsf T}A
= \tfrac12 S + \tfrac12 S^{\mathsf T} = S :
\]
故满射。$O_n = F^{-1}(I)$ 是一个光滑水平集，
维数为 $n^2 - \dim S_n = \frac{n(n-1)}2$。

\textbf{22.} 对 $A(t)^{\mathsf T}A(t) = I$ 求导，在 $t
= 0$ 处（$A(0) = I$）得 $A'(0)^{\mathsf T} + A'(0) = 0$，
即该矩阵反对称。反之，对 $K^{\mathsf T} = -K$，
$(\eu^{tK})^{\mathsf T}\eu^{tK} = \eu^{tK^{\mathsf T}}
\eu^{tK} = \eu^{-tK}\eu^{tK} = I$（对级数逐项\omterm{def:b2:linalg:transpose}{转置}，
指数可交换）：曲线留在 $O_n$ 中，
速度为 $K$，在 $t = 0$ 处。$I$ 处的切空间 $=$
反对称矩阵全体，其维数为预期的
$\frac{n(n-1)}2$。

\textbf{23.} $\operatorname{tr}K = 0$（对反对称矩阵 $K$），
故 $\det\eu^K = \eu^0 = 1$（第 18 题）：指数映射
落在 $SO_n$ 中。取 $n = 2$：$J^2 = -I$，故 $J^{2m} =
(-1)^mI$，$J^{2m+1} = (-1)^mJ$，且
\[
\eu^{\theta J}
= \Bigl(\sum_m\frac{(-1)^m\theta^{2m}}{(2m)!}\Bigr)I
+ \Bigl(\sum_m\frac{(-1)^m\theta^{2m+1}}{(2m+1)!}\Bigr)J
= \cos\theta\,I + \sin\theta\,J ,
\]
即旋转 $\theta$ 角：正弦与余弦级数
都包含在\omterm{ex:b2:nvs:matrixexp}{矩阵指数}之中。

\textbf{24.} 设 $A, B$ 可逆，则 $\operatorname{adj}(AB) =
\det(AB)(AB)^{-1} = \det(B)\det(A)B^{-1}A^{-1} =
\operatorname{adj}(B)\operatorname{adj}(A)$。等式两边
都是多项式（因而\omterm{def:b2:metric:continuity}{连续}）映射，其自变量为
$(A, B)$ 的元素；它们在稠密子集 $GL_n\times GL_n$
（位于 $\mathcal M_n\times\mathcal M_n$ 中，由第 3 题逼近
每个因子），故处处相等。

\textbf{25.} (i) 第一部分依赖于有限维赋范空间的\omterm{def:b2:metric:complete}{完备}性
（\omterm{def:b2:series:def}{绝对收敛}的级数收敛），第四部分
在此基础上还用到第 20 题的谱定理，第五部分的
局部对数依赖于反函数定理。(ii) \omterm{def:b2:diffcalc:differential}{微分}方程一章
依赖于刘维尔公式
（第 13 题）用于线性方程组的朗斯基\omterm{def:b2:linalg:det}{行列式}，并依赖于
$\frac{\dd}{\dd t}\eu^{tA} = A\eu^{tA}$（第 17 题），
这正是说 $\eu^{tA}$ 是 $X' = AX$ 的解。(iii)
$\dd(\det)_I = \operatorname{tr}$ 说明迹是
体积变化的无穷小速率，而 $\det\eu^A =
\eu^{\operatorname{tr}A}$ 则将该陈述整体积分。
(iv) 本科三年级分册为这些结构命名：$O_n$ 与
$SL_n(\R)$ 都是李群，它们在 $I$ 处的切空间
（反对称矩阵与迹为零的矩阵）是李代数，而
$\exp$ 是连接二者的桥梁。
\end{solution}
